\documentclass[12pt,a4paper]{article}

% Русский язык для pdfLaTeX
\usepackage[T2A]{fontenc}
\usepackage[utf8]{inputenc}
\usepackage[russian]{babel}

% Шрифт, близкий к Times New Roman.
% Если tempora не установлен, замени на mathptmx.
% \usepackage{tempora}
\usepackage{mathptmx}

\usepackage{amsmath,amssymb,amsthm}
\usepackage{mathtools}
\usepackage{xcolor}
\usepackage{listings}
\usepackage{enumitem}
\usepackage{tikz}
\usepackage{array}
\usepackage{booktabs}

% hyperref лучше подключать последним
\usepackage[unicode]{hyperref}

\newtheorem{definition}{Определение}
\newtheorem{theorem}{Теорема}
\newtheorem{lemma}{Лемма}
\newtheorem{property}{Свойство}

\lstset{
    basicstyle=\ttfamily\small,
    keywordstyle=\color{blue},
    commentstyle=\color{gray},
    stringstyle=\color{teal},
    frame=single,
    breaklines=true,
    columns=fullflexible,
    keepspaces=true,
    showstringspaces=false,
    extendedchars=false
}

\title{Билет №34 \\ \small Распределенное программирование. Процессы и их синхронизация. Объектно-ориентированное распределенное программирование. Параллельное программирование над общей памятью. (ИТМО) \\ Особенности программирования параллельных и распределенных систем. (СпБПУ)}
\author{Сериков Владислав Александрович}
\date{13 мая 2026}

\begin{document}

\maketitle
\tableofcontents
\newpage

\section{Базовые понятия параллельных и распределенных систем}

\begin{definition}
\textbf{Параллельная система} --- вычислительная система, в которой несколько вычислительных активностей выполняются одновременно или квазипараллельно для достижения общей цели.
\end{definition}

\begin{definition}
\textbf{Распределенная система} --- совокупность независимых вычислительных узлов, взаимодействующих посредством передачи сообщений и воспринимаемых пользователем или программой как единая система.
\end{definition}

Главное различие между параллельной и распределенной системой состоит не в количестве процессоров, а в модели взаимодействия:

\begin{center}
\begin{tabular}{p{0.25\textwidth}p{0.32\textwidth}p{0.32\textwidth}}
\toprule
\textbf{Критерий} & \textbf{Параллельная система} & \textbf{Распределенная система} \\
\midrule
Память & Часто общая память & Обычно раздельная память \\
Взаимодействие & Через общие переменные, блокировки, атомарные операции & Через сообщения, RPC, очереди, события \\
Отказы & Часто рассматриваются как отказ всей программы & Возможны частичные отказы узлов и каналов \\
Время & Иногда можно использовать общий системный таймер & Глобального надежного времени обычно нет \\
Цель & Ускорение вычислений & Масштабируемость, отказоустойчивость, географическое распределение \\
\bottomrule
\end{tabular}
\end{center}

\subsection{Процесс, поток, задача}

\begin{definition}
\textbf{Процесс} --- экземпляр выполняющейся программы, имеющий собственное адресное пространство, ресурсы ОС и поток управления.
\end{definition}

\begin{definition}
\textbf{Поток} --- легковесная единица выполнения внутри процесса, разделяющая с другими потоками адресное пространство процесса.
\end{definition}

\begin{definition}
\textbf{Задача} --- логическая единица работы. Задача может быть реализована процессом, потоком, корутиной, актором, удаленным вызовом или элементом очереди.
\end{definition}

В распределенном программировании обычно говорят о процессах как о независимых участниках вычисления. В параллельном программировании над общей памятью чаще говорят о потоках.

\section{Модели взаимодействия}

\subsection{Общая память}

В модели общей памяти несколько потоков имеют доступ к одним и тем же переменным.

\[
T_1, T_2, \dots, T_n \longrightarrow \text{общая память}
\]

Преимущества:
\begin{itemize}
    \item простота передачи данных: достаточно записать значение в общую переменную;
    \item высокая производительность на многоядерных системах;
    \item удобна для параллельных циклов, конвейеров, задач с разделяемыми структурами данных.
\end{itemize}

Недостатки:
\begin{itemize}
    \item гонки данных;
    \item взаимные блокировки;
    \item сложность доказательства корректности;
    \item проблемы кэш-когерентности и ложного разделения.
\end{itemize}

\subsection{Передача сообщений}

В модели передачи сообщений процессы не имеют общей памяти. Данные передаются явно.

\[
P_i \xrightarrow{\text{message}} P_j
\]

Операции:
\begin{itemize}
    \item \texttt{send(destination, message)};
    \item \texttt{receive(source, message)}.
\end{itemize}

Сообщения бывают:
\begin{itemize}
    \item синхронные и асинхронные;
    \item блокирующие и неблокирующие;
    \item точка-точка и коллективные;
    \item надежные и ненадежные;
    \item упорядоченные и неупорядоченные.
\end{itemize}

\subsection{Удаленный вызов процедур и методов}

\begin{definition}
\textbf{RPC} --- remote procedure call, механизм, позволяющий вызвать процедуру в другом адресном пространстве так, как если бы она была локальной.
\end{definition}

\begin{definition}
\textbf{RMI} --- remote method invocation, объектно-ориентированная разновидность RPC, при которой удаленно вызываются методы объектов.
\end{definition}

Типичная схема RPC/RMI:

\[
\text{client} \rightarrow \text{client stub} \rightarrow \text{network}
\rightarrow \text{server stub} \rightarrow \text{server object}
\]

Основные шаги:
\begin{enumerate}
    \item Клиент вызывает локальную функцию-заглушку.
    \item Заглушка сериализует параметры.
    \item Сообщение отправляется серверу.
    \item Серверная заглушка десериализует параметры.
    \item Вызывается реальная процедура или метод.
    \item Результат сериализуется и отправляется обратно.
    \item Клиентская заглушка возвращает результат клиенту.
\end{enumerate}

\section{Особенности программирования параллельных и распределенных систем}

\subsection{Недетерминизм}

Параллельная программа может иметь разные допустимые исполнения из-за различных расписаний потоков.

Например:

\begin{lstlisting}[language=C]
int x = 0;

Thread 1: x = x + 1;
Thread 2: x = x + 1;
\end{lstlisting}

Ожидаемый результат: $x = 2$.

Но операция \texttt{x = x + 1} не атомарна:
\[
\text{read } x,\quad \text{add } 1,\quad \text{write } x.
\]

Возможное чередование:

\begin{center}
\begin{tabular}{cccc}
\toprule
Шаг & Поток 1 & Поток 2 & $x$ \\
\midrule
1 & read $x=0$ & & 0 \\
2 & & read $x=0$ & 0 \\
3 & compute $1$ & & 0 \\
4 & & compute $1$ & 0 \\
5 & write $1$ & & 1 \\
6 & & write $1$ & 1 \\
\bottomrule
\end{tabular}
\end{center}

Итог: $x = 1$.

\subsection{Гонки данных}

\begin{definition}
\textbf{Гонка данных} возникает, когда два или более потока одновременно обращаются к одной ячейке памяти, хотя бы один доступ является записью, и между доступами нет синхронизации.
\end{definition}

Для устранения гонок применяются:
\begin{itemize}
    \item мьютексы;
    \item семафоры;
    \item мониторы;
    \item атомарные операции;
    \item барьеры;
    \item транзакционная память;
    \item неизменяемые структуры данных;
    \item передача сообщений вместо общей памяти.
\end{itemize}

\subsection{Взаимная блокировка}

\begin{definition}
\textbf{Deadlock}, или взаимная блокировка, --- ситуация, при которой группа процессов бесконечно ожидает событий, которые могут быть вызваны только процессами из этой же группы.
\end{definition}

Классические условия Коффмана:
\begin{enumerate}
    \item взаимное исключение;
    \item удержание и ожидание;
    \item отсутствие принудительного отъема ресурса;
    \item циклическое ожидание.
\end{enumerate}

Для предотвращения deadlock достаточно нарушить хотя бы одно из этих условий.

\subsection{Голодание и livelock}

\begin{definition}
\textbf{Голодание} --- ситуация, когда процесс теоретически может продолжить выполнение, но практически бесконечно не получает нужный ресурс.
\end{definition}

\begin{definition}
\textbf{Livelock} --- ситуация, когда процессы не заблокированы и продолжают выполнять действия, но полезного прогресса не происходит.
\end{definition}

\section{Синхронизация процессов и потоков}

\subsection{Критическая секция}

\begin{definition}
\textbf{Критическая секция} --- участок программы, в котором осуществляется доступ к разделяемому ресурсу и который должен выполняться не более чем одним процессом одновременно.
\end{definition}

Задача критической секции требует выполнения трех условий:
\begin{enumerate}
    \item \textbf{взаимное исключение}: в критической секции находится не более одного процесса;
    \item \textbf{прогресс}: если критическая секция свободна, выбор следующего процесса не откладывается бесконечно;
    \item \textbf{ограниченное ожидание}: процесс не должен ждать бесконечно долго.
\end{enumerate}

\subsection{Мьютекс}

\begin{definition}
\textbf{Мьютекс} --- объект синхронизации, предоставляющий взаимоисключающий доступ к ресурсу.
\end{definition}

Минимальный пример:

\begin{lstlisting}[language=C++]
std::mutex m;
int counter = 0;

void increment() {
    std::lock_guard<std::mutex> lock(m);
    counter++;
}
\end{lstlisting}

\subsection{Семафор}

\begin{definition}
\textbf{Семафор} --- целочисленный объект синхронизации с двумя атомарными операциями:
\[
P(S) \quad \text{и} \quad V(S).
\]
\end{definition}

Операция $P(S)$:
\begin{enumerate}
    \item если $S > 0$, уменьшить $S$ на $1$;
    \item иначе заблокировать процесс.
\end{enumerate}

Операция $V(S)$:
\begin{enumerate}
    \item увеличить $S$ на $1$;
    \item если есть ожидающие процессы, разбудить один из них.
\end{enumerate}

Пример ограничения числа одновременно работающих потоков:

\begin{lstlisting}[language=C++]
std::counting_semaphore<3> sem(3);

void worker() {
    sem.acquire();
    // не более трех потоков одновременно
    do_work();
    sem.release();
}
\end{lstlisting}

\subsection{Монитор}

\begin{definition}
\textbf{Монитор} --- высокоуровневая конструкция синхронизации, объединяющая данные, операции над ними и неявный механизм взаимного исключения.
\end{definition}

Монитор обычно содержит:
\begin{itemize}
    \item разделяемое состояние;
    \item методы доступа;
    \item условные переменные;
    \item операции \texttt{wait} и \texttt{signal}.
\end{itemize}

Пример схемы производителя и потребителя:

\begin{lstlisting}[language=Java]
class Buffer {
    private Queue<Integer> q = new LinkedList<>();
    private int capacity = 10;

    public synchronized void put(int x) throws InterruptedException {
        while (q.size() == capacity) {
            wait();
        }
        q.add(x);
        notifyAll();
    }

    public synchronized int get() throws InterruptedException {
        while (q.isEmpty()) {
            wait();
        }
        int x = q.remove();
        notifyAll();
        return x;
    }
}
\end{lstlisting}

\subsection{Барьер}

\begin{definition}
\textbf{Барьер} --- точка синхронизации, в которой каждый поток ожидает, пока остальные потоки не достигнут этой же точки.
\end{definition}

Схема:

\[
T_1,T_2,\dots,T_n \longrightarrow \text{barrier}
\longrightarrow \text{следующая фаза}
\]

Пример:

\begin{lstlisting}[language=C]
#pragma omp parallel
{
    phase1();
    #pragma omp barrier
    phase2();
}
\end{lstlisting}

\section{Классические алгоритмы синхронизации}

\subsection{Алгоритм Петерсона для двух процессов}

Алгоритм Петерсона решает задачу взаимного исключения для двух процессов с помощью двух флагов и переменной очередности.

Пусть:
\[
flag[i] = true
\]
означает, что процесс $i$ хочет войти в критическую секцию.

Переменная:
\[
turn \in \{0,1\}
\]
задает, кому процесс уступает право входа.

\subsubsection{Псевдокод}

Для процесса $i$, где $j = 1 - i$:

\begin{lstlisting}
flag[i] = true;
turn = j;

while (flag[j] && turn == j) {
    // ожидание
}

// critical section

flag[i] = false;

// remainder section
\end{lstlisting}

\subsubsection{Иллюстрация}

\begin{center}
\begin{tikzpicture}[node distance=2.5cm]
\node[draw, rectangle] (p0) {$P_0$: хочет войти};
\node[draw, rectangle, right of=p0, xshift=3cm] (p1) {$P_1$: хочет войти};
\node[draw, rectangle, below of=p0, xshift=3cm] (turn) {$turn$ указывает, кто уступает};
\node[draw, rectangle, below of=turn] (cs) {В критическую секцию входит один процесс};

\draw[->] (p0) -- (turn);
\draw[->] (p1) -- (turn);
\draw[->] (turn) -- (cs);
\end{tikzpicture}
\end{center}

\begin{theorem}
Алгоритм Петерсона обеспечивает взаимное исключение для двух процессов.
\end{theorem}

\begin{proof}
Предположим противное: оба процесса $P_0$ и $P_1$ одновременно находятся в критической секции.

Чтобы $P_0$ вошел в критическую секцию, условие его ожидания должно быть ложным:
\[
\neg(flag[1] \land turn = 1).
\]
Так как $P_1$ находится в критической секции, он перед входом установил:
\[
flag[1] = true.
\]
Следовательно, для входа $P_0$ необходимо:
\[
turn \neq 1,
\]
то есть:
\[
turn = 0.
\]

Аналогично, чтобы $P_1$ вошел в критическую секцию, условие его ожидания должно быть ложным:
\[
\neg(flag[0] \land turn = 0).
\]
Так как $P_0$ находится в критической секции, имеем:
\[
flag[0] = true.
\]
Следовательно, для входа $P_1$ необходимо:
\[
turn \neq 0,
\]
то есть:
\[
turn = 1.
\]

Получили одновременно:
\[
turn = 0 \quad \text{и} \quad turn = 1,
\]
что невозможно. Противоречие. Следовательно, оба процесса не могут находиться в критической секции одновременно.
\end{proof}

\begin{theorem}
Алгоритм Петерсона обеспечивает прогресс для двух процессов.
\end{theorem}

\begin{proof}
Если только один процесс, например $P_0$, хочет войти в критическую секцию, то:
\[
flag[1] = false.
\]
Тогда условие ожидания процесса $P_0$:
\[
flag[1] \land turn = 1
\]
ложно. Поэтому $P_0$ входит в критическую секцию.

Если оба процесса хотят войти, то они устанавливают свои флаги в $true$ и записывают переменную $turn$. Последняя запись в $turn$ определяет, кто уступает. Если последней была запись $turn = 1$, то $P_0$ ждет, а $P_1$ проходит. Если последней была запись $turn = 0$, то $P_1$ ждет, а $P_0$ проходит.

Следовательно, когда критическая секция свободна, хотя бы один желающий процесс сможет в нее войти.
\end{proof}

\subsection{Задача производителя и потребителя}

Есть буфер ограниченного размера. Производители помещают элементы, потребители извлекают.

Используются три семафора:
\[
empty = N,\quad full = 0,\quad mutex = 1.
\]

\subsubsection{Производитель}

\begin{lstlisting}
while true:
    item = produce()
    P(empty)
    P(mutex)
    buffer.add(item)
    V(mutex)
    V(full)
\end{lstlisting}

\subsubsection{Потребитель}

\begin{lstlisting}
while true:
    P(full)
    P(mutex)
    item = buffer.remove()
    V(mutex)
    V(empty)
    consume(item)
\end{lstlisting}

\subsubsection{Смысл семафоров}

\begin{itemize}
    \item $empty$ считает свободные ячейки;
    \item $full$ считает занятые ячейки;
    \item $mutex$ защищает сам буфер.
\end{itemize}

\section{Параллельное программирование над общей памятью}

\subsection{Модель fork-join}

В модели fork-join программа начинается одним управляющим потоком. Затем создается группа потоков, выполняющих параллельную область, после чего они синхронизируются и снова остается один поток.

\[
\text{один поток} \rightarrow
\text{параллельная область} \rightarrow
\text{один поток}
\]

Иллюстрация:

\begin{center}
\begin{tikzpicture}
\node at (0,0) (a) {master};
\draw[->] (0.8,0) -- (2,0);

\draw[->] (2,0) -- (4,1);
\draw[->] (2,0) -- (4,0.3);
\draw[->] (2,0) -- (4,-0.3);
\draw[->] (2,0) -- (4,-1);

\node at (4.7,1) {$T_1$};
\node at (4.7,0.3) {$T_2$};
\node at (4.7,-0.3) {$T_3$};
\node at (4.7,-1) {$T_4$};

\draw[->] (5.2,1) -- (7,0);
\draw[->] (5.2,0.3) -- (7,0);
\draw[->] (5.2,-0.3) -- (7,0);
\draw[->] (5.2,-1) -- (7,0);

\node at (8,0) {master};
\end{tikzpicture}
\end{center}

\subsection{OpenMP как модель программирования общей памяти}

OpenMP предоставляет:
\begin{itemize}
    \item директивы компилятора;
    \item функции времени выполнения;
    \item переменные окружения.
\end{itemize}

Пример параллельного цикла:

\begin{lstlisting}[language=C]
#pragma omp parallel for
for (int i = 0; i < n; i++) {
    a[i] = b[i] + c[i];
}
\end{lstlisting}

Пример редукции:

\begin{lstlisting}[language=C]
double sum = 0.0;

#pragma omp parallel for reduction(+:sum)
for (int i = 0; i < n; i++) {
    sum += a[i];
}
\end{lstlisting}

Без редукции здесь возникла бы гонка данных по переменной \texttt{sum}.

\subsection{Атомарные операции}

\begin{definition}
\textbf{Атомарная операция} --- операция, которая выполняется неделимо относительно других потоков.
\end{definition}

Пример:

\begin{lstlisting}[language=C++]
std::atomic<int> counter = 0;

void increment() {
    counter.fetch_add(1);
}
\end{lstlisting}

Атомарные операции часто реализуются аппаратными инструкциями:
\begin{itemize}
    \item test-and-set;
    \item compare-and-swap;
    \item fetch-and-add;
    \item load-linked/store-conditional.
\end{itemize}

\subsection{Compare-and-swap}

Операция CAS имеет вид:

\[
CAS(addr, expected, desired).
\]

Она атомарно:
\begin{enumerate}
    \item читает значение по адресу $addr$;
    \item сравнивает его с $expected$;
    \item если значения равны, записывает $desired$;
    \item возвращает успех или неуспех операции.
\end{enumerate}

Пример реализации неблокирующего инкремента:

\begin{lstlisting}[language=C++]
std::atomic<int> x = 0;

void increment() {
    int old = x.load();
    while (!x.compare_exchange_weak(old, old + 1)) {
        // old автоматически обновляется текущим значением x
    }
}
\end{lstlisting}

\section{Корректность параллельных объектов}

\subsection{Линеаризуемость}

\begin{definition}
\textbf{Линеаризуемость} --- свойство параллельного объекта, означающее, что каждая операция выглядит как мгновенно выполненная в некоторой точке между своим вызовом и возвратом, а получившийся последовательный порядок удовлетворяет спецификации объекта.
\end{definition}

Такая точка называется \textbf{точкой линеаризации}.

Пример: потокобезопасный стек.

\begin{lstlisting}
Thread 1: push(1) starts
Thread 2: pop() starts
Thread 1: push(1) returns
Thread 2: pop() returns 1
\end{lstlisting}

Такое исполнение линеаризуемо, если можно считать, что \texttt{push(1)} произошел до \texttt{pop()}.

\begin{theorem}
Если каждый объект в системе линеаризуем, то вся система объектов линеаризуема при композиции.
\end{theorem}

\begin{proof}
Пусть в системе есть объекты:
\[
O_1, O_2, \dots, O_k.
\]
Для каждого объекта $O_i$ существует корректный последовательный порядок операций $S_i$, который:
\begin{enumerate}
    \item согласован с реальным временем для неперекрывающихся операций над $O_i$;
    \item удовлетворяет последовательной спецификации $O_i$.
\end{enumerate}

Рассмотрим объединение всех операций над всеми объектами. Для каждой операции выберем ее точку линеаризации внутри интервала между вызовом и ответом. Так как для каждого объекта такие точки существуют, можно упорядочить все операции глобально по этим точкам.

Если две операции относятся к одному объекту, их порядок совпадает с $S_i$, следовательно, спецификация этого объекта не нарушается. Если операции относятся к разным объектам, между ними нет общей объектной спецификации, кроме требования согласованности с реальным временем. Поскольку точки линеаризации лежат внутри интервалов операций, если операция $a$ завершилась до начала операции $b$, то точка линеаризации $a$ обязательно предшествует точке линеаризации $b$.

Следовательно, глобальный порядок всех операций удовлетворяет спецификациям всех объектов и сохраняет реальное время. Значит, композиция линеаризуемых объектов линеаризуема.
\end{proof}

\section{Распределенное программирование}

\subsection{Основные свойства распределенных систем}

Распределенные системы характеризуются следующими особенностями:
\begin{itemize}
    \item отсутствие общей памяти;
    \item отсутствие общего надежного времени;
    \item задержки передачи сообщений;
    \item частичные отказы;
    \item независимое выполнение процессов;
    \item необходимость согласования состояния;
    \item сложность наблюдения глобального состояния.
\end{itemize}

\subsection{Модель распределенного вычисления}

Пусть есть множество процессов:
\[
P = \{P_1, P_2, \dots, P_n\}.
\]

Каждый процесс выполняет последовательность событий:
\[
e_{i1}, e_{i2}, \dots
\]

События бывают:
\begin{itemize}
    \item внутренние;
    \item отправка сообщения;
    \item получение сообщения.
\end{itemize}

\subsection{Отношение ``произошло раньше''}

\begin{definition}
Отношение \textbf{happened-before}, обозначаемое $a \rightarrow b$, определяется так:
\begin{enumerate}
    \item если $a$ и $b$ --- события одного процесса и $a$ произошло раньше $b$, то $a \rightarrow b$;
    \item если $a$ --- отправка сообщения, а $b$ --- получение этого сообщения, то $a \rightarrow b$;
    \item если $a \rightarrow b$ и $b \rightarrow c$, то $a \rightarrow c$.
\end{enumerate}
\end{definition}

Если ни $a \rightarrow b$, ни $b \rightarrow a$, то события называются \textbf{конкурентными}:
\[
a \parallel b.
\]

\subsection{Логические часы Лампорта}

\begin{definition}
\textbf{Логические часы Лампорта} --- механизм присваивания событиям числовых меток времени так, чтобы из причинной зависимости следовало возрастание меток.
\end{definition}

Каждый процесс $P_i$ хранит счетчик $C_i$.

\subsubsection{Алгоритм}

\begin{enumerate}
    \item Перед каждым локальным событием процесс увеличивает счетчик:
    \[
    C_i := C_i + 1.
    \]
    \item При отправке сообщения $m$ процесс прикрепляет к нему метку:
    \[
    timestamp(m) = C_i.
    \]
    \item При получении сообщения $m$ с меткой $t$ процесс выполняет:
    \[
    C_j := \max(C_j, t) + 1.
    \]
\end{enumerate}

\subsubsection{Минимальный пример}

Пусть:

\[
P_1: a \rightarrow b,\qquad P_2: c \rightarrow d.
\]

Событие $b$ --- отправка сообщения, событие $d$ --- получение.

\begin{center}
\begin{tikzpicture}
\node at (0,2) {$P_1$};
\node at (0,0) {$P_2$};

\draw[->] (1,2) -- (7,2);
\draw[->] (1,0) -- (7,0);

\node[circle,draw] (a) at (2,2) {$a:1$};
\node[circle,draw] (b) at (4,2) {$b:2$};

\node[circle,draw] (c) at (2,0) {$c:1$};
\node[circle,draw] (d) at (6,0) {$d:3$};

\draw[->, dashed] (b) -- node[right] {msg, $t=2$} (d);
\end{tikzpicture}
\end{center}

При получении сообщения:
\[
C_2 := \max(1,2)+1 = 3.
\]

\begin{theorem}
Если $a \rightarrow b$, то для логических часов Лампорта выполняется:
\[
C(a) < C(b).
\]
\end{theorem}

\begin{proof}
Докажем по построению отношения $\rightarrow$.

\textbf{Случай 1.} $a$ и $b$ --- события одного процесса, и $a$ произошло раньше $b$.

По правилу алгоритма счетчик процесса увеличивается перед каждым событием. Следовательно, более позднее событие получит строго большую метку:
\[
C(a) < C(b).
\]

\textbf{Случай 2.} $a$ --- отправка сообщения, $b$ --- получение этого сообщения.

Пусть сообщение отправлено с меткой:
\[
t = C(a).
\]
При получении процесс устанавливает:
\[
C(b) = \max(C_{\text{local}}, t) + 1.
\]
Следовательно:
\[
C(b) > t = C(a).
\]

\textbf{Случай 3.} $a \rightarrow b$ следует из транзитивности: существуют события
\[
a = e_0 \rightarrow e_1 \rightarrow \dots \rightarrow e_k = b.
\]
По первым двум случаям для каждой пары:
\[
C(e_r) < C(e_{r+1}).
\]
Следовательно:
\[
C(a) = C(e_0) < C(e_1) < \dots < C(e_k) = C(b).
\]

Во всех случаях $C(a) < C(b)$.
\end{proof}

Обратное утверждение неверно:
\[
C(a) < C(b) \not\Rightarrow a \rightarrow b.
\]
Разные независимые события могут случайно получить упорядоченные метки.

\subsection{Векторные часы}

\begin{definition}
\textbf{Векторные часы} --- механизм логического времени, в котором каждому событию сопоставляется вектор длины $n$, позволяющий определять причинную зависимость точнее, чем часы Лампорта.
\end{definition}

Каждый процесс $P_i$ хранит вектор:
\[
V_i[1..n].
\]

\subsubsection{Алгоритм}

\begin{enumerate}
    \item Перед локальным событием процесс $P_i$ выполняет:
    \[
    V_i[i] := V_i[i] + 1.
    \]
    \item При отправке сообщения процесс прикрепляет к нему копию вектора $V_i$.
    \item При получении сообщения с вектором $T$ процесс $P_j$ выполняет:
    \[
    V_j[k] := \max(V_j[k], T[k]) \quad \forall k,
    \]
    затем:
    \[
    V_j[j] := V_j[j] + 1.
    \]
\end{enumerate}

Порядок векторов определяется так:
\[
V < W
\]
если:
\[
\forall k: V[k] \leq W[k]
\]
и
\[
\exists r: V[r] < W[r].
\]

\begin{theorem}
Для векторных часов:
\[
a \rightarrow b \iff V(a) < V(b).
\]
\end{theorem}

\begin{proof}
Докажем обе стороны.

\textbf{1. Если $a \rightarrow b$, то $V(a) < V(b)$.}

Рассмотрим базовые случаи.

Если $a$ и $b$ --- события одного процесса, причем $a$ раньше $b$, то компонент этого процесса увеличивается при каждом событии. Остальные компоненты не уменьшаются. Поэтому:
\[
V(a) \leq V(b)
\]
покомпонентно, и хотя бы одна компонента строго возрастает. Значит:
\[
V(a) < V(b).
\]

Если $a$ --- отправка сообщения, а $b$ --- его получение, то вектор отправителя прикрепляется к сообщению. Получатель берет покомпонентный максимум со своим вектором, а затем увеличивает собственную компоненту. Следовательно, все компоненты $V(a)$ не превосходят соответствующих компонент $V(b)$, и хотя бы одна компонента строго больше. Значит:
\[
V(a) < V(b).
\]

Для транзитивного случая рассуждение аналогично: если
\[
a = e_0 \rightarrow e_1 \rightarrow \dots \rightarrow e_k = b,
\]
то
\[
V(e_0) < V(e_1) < \dots < V(e_k),
\]
откуда:
\[
V(a) < V(b).
\]

\textbf{2. Если $V(a) < V(b)$, то $a \rightarrow b$.}

Компонента $V(x)[i]$ показывает, сколько событий процесса $P_i$ причинно известно событию $x$. Если:
\[
V(a) < V(b),
\]
то все события, отраженные в $V(a)$, входят в причинное прошлое события $b$, и хотя бы одно знание у $b$ строго больше. В частности, событие $a$ учтено в причинной истории $b$. Это возможно только если информация о событии $a$ достигла процесса, на котором произошло $b$, по цепочке локальных шагов и сообщений. По определению happened-before такая цепочка означает:
\[
a \rightarrow b.
\]

Следовательно:
\[
a \rightarrow b \iff V(a) < V(b).
\]
\end{proof}

Если два вектора несравнимы, соответствующие события конкурентны.

\section{Глобальное состояние и снимки}

\subsection{Проблема глобального состояния}

В распределенной системе нельзя просто остановить все процессы одновременно, потому что:
\begin{itemize}
    \item нет общего времени;
    \item сообщения могут находиться в каналах;
    \item процессы выполняются независимо;
    \item остановка всей системы может быть невозможна.
\end{itemize}

\begin{definition}
\textbf{Глобальное состояние} распределенной системы состоит из локальных состояний всех процессов и состояний всех каналов связи.
\end{definition}

\begin{definition}
\textbf{Согласованный срез} --- множество событий, замкнутое относительно причинности: если событие принадлежит срезу, то все его причинные предшественники также принадлежат срезу.
\end{definition}

Формально, срез $C$ согласован, если:
\[
e \in C \land f \rightarrow e \Rightarrow f \in C.
\]

\subsection{Алгоритм Чанди--Лампорта}

Алгоритм Чанди--Лампорта строит согласованный глобальный снимок в асинхронной системе с надежными FIFO-каналами.

\subsubsection{Предположения}

\begin{enumerate}
    \item Процессы не отказывают.
    \item Каналы надежны.
    \item Каналы FIFO.
    \item Между каждой парой процессов есть направленные каналы.
    \item Специальные маркерные сообщения не теряются.
\end{enumerate}

\subsubsection{Алгоритм}

Инициатор:
\begin{enumerate}
    \item записывает свое локальное состояние;
    \item отправляет маркер по всем исходящим каналам;
    \item начинает записывать сообщения, приходящие по входящим каналам, пока по ним не придет маркер.
\end{enumerate}

Обычный процесс при первом получении маркера по каналу $c$:
\begin{enumerate}
    \item записывает свое локальное состояние;
    \item считает состояние канала $c$ пустым;
    \item отправляет маркер по всем исходящим каналам;
    \item начинает записывать сообщения, приходящие по другим входящим каналам, пока по ним не придет маркер.
\end{enumerate}

Если процесс уже записал свое состояние и получает маркер по каналу $c$:
\begin{enumerate}
    \item прекращает запись сообщений по каналу $c$;
    \item записанные до маркера сообщения считаются состоянием канала $c$.
\end{enumerate}

\subsubsection{Иллюстрация}

\begin{center}
\begin{tikzpicture}[node distance=4cm]
\node[circle,draw] (p1) {$P_1$};
\node[circle,draw,right of=p1] (p2) {$P_2$};
\node[circle,draw,below of=p1] (p3) {$P_3$};

\draw[->] (p1) -- node[above] {marker} (p2);
\draw[->] (p1) -- node[left] {marker} (p3);
\draw[->] (p2) -- (p3);
\draw[->] (p3) -- (p1);
\end{tikzpicture}
\end{center}

\begin{theorem}
Алгоритм Чанди--Лампорта строит согласованный глобальный снимок распределенной системы.
\end{theorem}

\begin{proof}
Нужно доказать, что построенный срез замкнут относительно причинности.

Пусть событие $e$ принадлежит снимку, а событие $f$ таково, что:
\[
f \rightarrow e.
\]
Нужно показать, что $f$ также принадлежит снимку.

Рассмотрим возможные случаи.

\textbf{1. События $f$ и $e$ произошли в одном процессе.}

Если $e$ принадлежит локальному состоянию процесса в момент записи, то все более ранние события этого процесса также уже произошли до записи локального состояния. Так как $f$ предшествует $e$ в том же процессе, $f$ также принадлежит снимку.

\textbf{2. $f$ --- отправка сообщения, а $e$ --- получение этого сообщения.}

Пусть сообщение $m$ отправлено в событии $f$ и получено в событии $e$.

Если получение $e$ вошло в локальное состояние получателя, значит оно произошло до записи состояния получателя.

Нужно показать, что отправка $f$ также отражена в снимке.

Предположим противное: отправка $f$ не вошла в снимок отправителя. Тогда отправитель записал свое состояние до отправки $m$. По алгоритму сразу после записи состояния отправитель посылает маркеры по всем исходящим каналам. Так как каналы FIFO, маркер по каналу от отправителя к получателю должен прийти раньше сообщения $m$, если $m$ было отправлено после маркера.

Но получатель записывает свое состояние при первом получении маркера. Если маркер пришел до $m$, то получение $m$ не могло произойти до записи состояния получателя. Это противоречит тому, что $e$ вошло в локальное состояние получателя.

Следовательно, отправка $f$ вошла в снимок.

\textbf{3. Транзитивный случай.}

Если:
\[
f = e_0 \rightarrow e_1 \rightarrow \dots \rightarrow e_k = e,
\]
то по первым двум случаям принадлежность $e_r$ снимку влечет принадлежность $e_{r-1}$ снимку. По индукции получаем, что $f$ принадлежит снимку.

Следовательно, снимок замкнут относительно happened-before и является согласованным.
\end{proof}

\section{Объектно-ориентированное распределенное программирование}

\subsection{Идея распределенных объектов}

В объектно-ориентированном распределенном программировании объект может находиться в одном процессе или на одном узле, а его методы вызываются из другого процесса или узла.

\begin{definition}
\textbf{Распределенный объект} --- объект, доступ к методам которого может осуществляться удаленно через сетевое взаимодействие.
\end{definition}

Основные элементы:
\begin{itemize}
    \item удаленный интерфейс;
    \item реализация объекта;
    \item прокси или stub на стороне клиента;
    \item skeleton или dispatcher на стороне сервера;
    \item механизм сериализации;
    \item реестр или служба имен.
\end{itemize}

\subsection{Отличия локальных и удаленных вызовов}

\begin{center}
\begin{tabular}{p{0.3\textwidth}p{0.3\textwidth}p{0.3\textwidth}}
\toprule
\textbf{Свойство} & \textbf{Локальный вызов} & \textbf{Удаленный вызов} \\
\midrule
Время выполнения & Наносекунды--микросекунды & Миллисекунды и более \\
Отказ & Обычно исключение процесса & Возможен отказ сети, сервера, тайм-аут \\
Передача параметров & По значению или ссылке & Обычно сериализация, копирование, удаленные ссылки \\
Идентичность объекта & Адрес в памяти & Глобальный идентификатор или удаленная ссылка \\
Побочные эффекты & Локальные & Распределенные, частично выполненные \\
\bottomrule
\end{tabular}
\end{center}

\subsection{Семантика удаленного вызова}

Возможные гарантии:
\begin{itemize}
    \item \textbf{maybe}: вызов мог выполниться, а мог не выполниться;
    \item \textbf{at-least-once}: вызов выполнится хотя бы один раз, но возможны повторы;
    \item \textbf{at-most-once}: вызов выполнится не более одного раза;
    \item \textbf{exactly-once}: идеальная семантика, трудно достижимая при отказах.
\end{itemize}

На практике часто реализуют \textbf{at-most-once} с помощью уникальных идентификаторов запросов и кэша ответов.

\subsection{Минимальный пример RMI-подобного интерфейса}

\begin{lstlisting}[language=Java]
public interface Calculator extends Remote {
    int add(int a, int b) throws RemoteException;
}
\end{lstlisting}

Реализация:

\begin{lstlisting}[language=Java]
public class CalculatorImpl implements Calculator {
    public int add(int a, int b) {
        return a + b;
    }
}
\end{lstlisting}

Клиент:

\begin{lstlisting}[language=Java]
Calculator calc = lookup("CalculatorService");
int result = calc.add(2, 3);
\end{lstlisting}

Логически клиент вызывает метод \texttt{add}, но физически происходит сериализация, сетевой обмен и выполнение метода на сервере.

\subsection{Акторная модель}

\begin{definition}
\textbf{Актор} --- независимый объект вычисления, имеющий собственное состояние и почтовый ящик, взаимодействующий с другими акторами только через асинхронные сообщения.
\end{definition}

Актор при получении сообщения может:
\begin{enumerate}
    \item изменить свое состояние;
    \item отправить сообщения другим акторам;
    \item создать новых акторов.
\end{enumerate}

Преимущества:
\begin{itemize}
    \item отсутствие общей изменяемой памяти между акторами;
    \item естественная масштабируемость;
    \item удобство моделирования распределенных систем;
    \item изоляция отказов.
\end{itemize}

\section{MPI и программирование с передачей сообщений}

\subsection{Общая идея MPI}

MPI --- стандарт программирования с передачей сообщений. Обычно используется для высокопроизводительных вычислений на кластерах.

Основные понятия:
\begin{itemize}
    \item \textbf{communicator} --- группа процессов;
    \item \textbf{rank} --- номер процесса внутри коммуникатора;
    \item \textbf{point-to-point communication} --- обмен между двумя процессами;
    \item \textbf{collective communication} --- коллективные операции над группой процессов.
\end{itemize}

\subsection{Минимальная MPI-программа}

\begin{lstlisting}[language=C]
#include <mpi.h>
#include <stdio.h>

int main(int argc, char** argv) {
    MPI_Init(&argc, &argv);

    int rank, size;
    MPI_Comm_rank(MPI_COMM_WORLD, &rank);
    MPI_Comm_size(MPI_COMM_WORLD, &size);

    printf("Hello from process %d of %d\n", rank, size);

    MPI_Finalize();
    return 0;
}
\end{lstlisting}

\subsection{Передача сообщения}

\begin{lstlisting}[language=C]
if (rank == 0) {
    int x = 42;
    MPI_Send(&x, 1, MPI_INT, 1, 0, MPI_COMM_WORLD);
} else if (rank == 1) {
    int y;
    MPI_Recv(&y, 1, MPI_INT, 0, 0, MPI_COMM_WORLD,
             MPI_STATUS_IGNORE);
    printf("%d\n", y);
}
\end{lstlisting}

\subsection{Коллективные операции}

Примеры:
\begin{itemize}
    \item \texttt{MPI\_Bcast} --- широковещательная рассылка;
    \item \texttt{MPI\_Scatter} --- распределение частей массива;
    \item \texttt{MPI\_Gather} --- сбор частей массива;
    \item \texttt{MPI\_Reduce} --- редукция;
    \item \texttt{MPI\_Barrier} --- барьер.
\end{itemize}

Пример суммы:

\begin{lstlisting}[language=C]
int local = rank + 1;
int global = 0;

MPI_Reduce(&local, &global, 1, MPI_INT, MPI_SUM, 0,
           MPI_COMM_WORLD);

if (rank == 0) {
    printf("sum = %d\n", global);
}
\end{lstlisting}

\section{Распределенная взаимная блокировка и ее обнаружение}

В распределенной системе deadlock сложнее обнаруживать, потому что нет единого места, где известно все состояние.

\subsection{Граф ожидания}

\begin{definition}
\textbf{Граф ожидания} --- ориентированный граф, вершинами которого являются процессы, а ребро
\[
P_i \rightarrow P_j
\]
означает, что процесс $P_i$ ожидает ресурс или сообщение от $P_j$.
\end{definition}

\begin{theorem}
В системе с ресурсами единичной мощности deadlock существует тогда и только тогда, когда в графе ожидания есть ориентированный цикл.
\end{theorem}

\begin{proof}
\textbf{Необходимость.}

Пусть есть deadlock. Рассмотрим множество процессов $D$, находящихся во взаимной блокировке. Каждый процесс из $D$ ожидает ресурс, удерживаемый некоторым другим процессом из $D$, иначе он мог бы продолжить выполнение. Значит, из каждой вершины $P_i \in D$ выходит ребро к некоторой вершине $P_j \in D$.

Так как множество $D$ конечно, если двигаться по исходящим ребрам, рано или поздно некоторая вершина повторится. Повторение вершины образует ориентированный цикл.

\textbf{Достаточность.}

Пусть в графе ожидания есть цикл:
\[
P_1 \rightarrow P_2 \rightarrow \dots \rightarrow P_k \rightarrow P_1.
\]
Это означает, что $P_1$ ждет ресурс у $P_2$, $P_2$ ждет ресурс у $P_3$, ..., $P_k$ ждет ресурс у $P_1$.

Так как ресурсы единичной мощности и удерживаются ожидающими процессами, ни один процесс из цикла не может продолжить выполнение, пока другой не освободит ресурс. Но освобождение невозможно без продолжения выполнения. Следовательно, процессы цикла находятся во взаимной блокировке.
\end{proof}

\section{Согласование и консенсус}

\subsection{Задача консенсуса}

\begin{definition}
\textbf{Консенсус} --- задача, в которой несколько процессов должны выбрать одно общее значение, удовлетворяющее условиям корректности.
\end{definition}

Условия:
\begin{enumerate}
    \item \textbf{Termination}: каждый корректный процесс рано или поздно принимает решение.
    \item \textbf{Agreement}: никакие два корректных процесса не принимают разные решения.
    \item \textbf{Validity}: принятое значение должно быть предложено некоторым процессом.
\end{enumerate}

Консенсус важен для:
\begin{itemize}
    \item репликации;
    \item распределенных транзакций;
    \item выбора лидера;
    \item журналов команд;
    \item отказоустойчивых баз данных.
\end{itemize}

\subsection{Невозможность FLP}

\begin{theorem}[FLP, неформальная формулировка]
В полностью асинхронной распределенной системе с передачей сообщений невозможно построить детерминированный алгоритм консенсуса, который гарантирует завершение при возможности отказа хотя бы одного процесса остановкой.
\end{theorem}

Идея доказательства:
\begin{itemize}
    \item в асинхронной системе невозможно отличить медленный процесс от отказавшего;
    \item существуют бивалентные конфигурации, из которых еще возможны разные решения;
    \item можно строить бесконечное исполнение, в котором система остается бивалентной;
    \item значит, гарантированного завершения нет.
\end{itemize}

Практические алгоритмы консенсуса обходят ограничение за счет:
\begin{itemize}
    \item тайм-аутов;
    \item частичной синхронности;
    \item случайности;
    \item предположений о большинстве;
    \item детекторов отказов.
\end{itemize}

\section{Транзакции и двухфазная фиксация}

\subsection{Распределенная транзакция}

\begin{definition}
\textbf{Распределенная транзакция} --- транзакция, затрагивающая несколько узлов или ресурсов.
\end{definition}

Свойства ACID:
\begin{itemize}
    \item Atomicity --- атомарность;
    \item Consistency --- согласованность;
    \item Isolation --- изоляция;
    \item Durability --- долговечность.
\end{itemize}

\subsection{Алгоритм двухфазной фиксации}

Участники:
\begin{itemize}
    \item координатор;
    \item участники транзакции.
\end{itemize}

\subsubsection{Фаза 1: prepare}

\begin{enumerate}
    \item Координатор отправляет всем участникам сообщение \texttt{PREPARE}.
    \item Каждый участник проверяет, может ли зафиксировать транзакцию.
    \item Если может, записывает состояние \texttt{prepared} в журнал и отвечает \texttt{VOTE\_COMMIT}.
    \item Иначе отвечает \texttt{VOTE\_ABORT}.
\end{enumerate}

\subsubsection{Фаза 2: commit/abort}

\begin{enumerate}
    \item Если все участники ответили \texttt{VOTE\_COMMIT}, координатор записывает \texttt{GLOBAL\_COMMIT} и рассылает \texttt{COMMIT}.
    \item Если хотя бы один участник ответил \texttt{VOTE\_ABORT}, координатор записывает \texttt{GLOBAL\_ABORT} и рассылает \texttt{ABORT}.
    \item Участники выполняют решение и подтверждают его.
\end{enumerate}

Недостаток: 2PC является блокирующим протоколом. Если координатор отказал после того, как участник проголосовал за commit, но до получения глобального решения, участник может быть вынужден ждать восстановления координатора.

\section{Проблемы производительности}

\subsection{Закон Амдала}

\begin{theorem}[Закон Амдала]
Если доля последовательной части программы равна $s$, а остальная часть идеально распараллеливается на $p$ процессоров, то ускорение равно:
\[
S(p) = \frac{1}{s + \frac{1-s}{p}}.
\]
\end{theorem}

\begin{proof}
Пусть исходное время выполнения программы на одном процессоре равно $T$.

Последовательная часть занимает:
\[
sT.
\]

Параллельная часть занимает:
\[
(1-s)T.
\]

При идеальном распараллеливании на $p$ процессоров параллельная часть выполнится за:
\[
\frac{(1-s)T}{p}.
\]

Итоговое время:
\[
T_p = sT + \frac{(1-s)T}{p}
= T\left(s + \frac{1-s}{p}\right).
\]

Ускорение:
\[
S(p) = \frac{T}{T_p}
= \frac{T}{T\left(s + \frac{1-s}{p}\right)}
= \frac{1}{s + \frac{1-s}{p}}.
\]
\end{proof}

Следствие:
\[
\lim_{p \to \infty} S(p) = \frac{1}{s}.
\]

Даже малая последовательная часть ограничивает максимальное ускорение.

\subsection{Закон Густафсона}

Если при увеличении числа процессоров увеличивать размер задачи, то ускорение можно оценить как:
\[
S_G(p) = p - s(p-1).
\]

В отличие от закона Амдала, закон Густафсона показывает, почему большие задачи могут хорошо масштабироваться.

\subsection{Накладные расходы параллелизма}

Основные источники накладных расходов:
\begin{itemize}
    \item создание потоков и процессов;
    \item синхронизация;
    \item передача сообщений;
    \item ожидание на барьерах;
    \item дисбаланс нагрузки;
    \item промахи кэша;
    \item ложное разделение;
    \item сериализация и десериализация;
    \item повторные вычисления после отказов.
\end{itemize}

\subsection{Ложное разделение}

\begin{definition}
\textbf{Ложное разделение} --- ситуация, когда разные потоки изменяют разные переменные, находящиеся в одной кэш-линии, из-за чего кэш-линия постоянно инвалидируется.
\end{definition}

Пример:

\begin{lstlisting}[language=C++]
struct Counters {
    std::atomic<long> a;
    std::atomic<long> b;
};
\end{lstlisting}

Если \texttt{a} и \texttt{b} лежат в одной кэш-линии, два потока могут мешать друг другу, хотя логически не разделяют одну переменную.

\section{Типичные архитектурные подходы}

\subsection{Master--worker}

Один главный процесс распределяет задачи между рабочими.

Плюсы:
\begin{itemize}
    \item простота реализации;
    \item централизованный контроль;
    \item удобно для независимых задач.
\end{itemize}

Минусы:
\begin{itemize}
    \item master может стать узким местом;
    \item отказ master критичен;
    \item возможен дисбаланс нагрузки.
\end{itemize}

\subsection{Pipeline}

Вычисление разбивается на стадии:

\[
S_1 \rightarrow S_2 \rightarrow \dots \rightarrow S_k.
\]

Каждая стадия обрабатывает данные и передает результат следующей.

\subsection{MapReduce}

Общая схема:
\begin{enumerate}
    \item \textbf{Map}: преобразование входных данных в пары ключ--значение.
    \item \textbf{Shuffle}: группировка по ключам.
    \item \textbf{Reduce}: агрегация значений одного ключа.
\end{enumerate}

Пример подсчета слов:

\[
map(\text{``a b a''}) = (a,1),(b,1),(a,1)
\]

\[
reduce(a,[1,1]) = (a,2)
\]

\[
reduce(b,[1]) = (b,1)
\]

\section{Отказоустойчивость}

\subsection{Типы отказов}

\begin{itemize}
    \item crash failure --- процесс останавливается;
    \item omission failure --- сообщение не отправлено или не получено;
    \item timing failure --- нарушение временных ограничений;
    \item Byzantine failure --- произвольное, в том числе злонамеренное поведение.
\end{itemize}

\subsection{Репликация}

\begin{definition}
\textbf{Репликация} --- хранение копий данных или сервисов на нескольких узлах.
\end{definition}

Виды:
\begin{itemize}
    \item активная репликация;
    \item пассивная репликация;
    \item primary-backup;
    \item quorum-based replication.
\end{itemize}

\subsection{Кворумы}

Пусть:
\[
N
\]
--- число реплик.

Для операций чтения и записи выбираются кворумы $R$ и $W$.

Чтобы чтение пересекалось с последней записью, требуется:
\[
R + W > N.
\]

Чтобы две записи пересекались:
\[
W > \frac{N}{2}.
\]

\section{CAP-теорема}

\begin{theorem}[CAP, формулировка в инженерном смысле]
В распределенной системе хранения данных при сетевом разделении невозможно одновременно гарантировать строгую согласованность и доступность для всех запросов.
\end{theorem}

Здесь:
\begin{itemize}
    \item \textbf{Consistency} --- все клиенты видят согласованное состояние;
    \item \textbf{Availability} --- каждый запрос к неотказавшему узлу получает ответ;
    \item \textbf{Partition tolerance} --- система продолжает работать при разрывах сети.
\end{itemize}

\begin{proof}
Рассмотрим систему с двумя репликами $A$ и $B$. Между ними возникло сетевое разделение: сообщения от $A$ к $B$ и от $B$ к $A$ не доставляются.

Клиент $C_1$ отправляет запись:
\[
x := 1
\]
на реплику $A$.

Если система должна быть доступной, $A$ обязана ответить клиенту, не ожидая связи с $B$.

Теперь клиент $C_2$ читает $x$ с реплики $B$.

Так как разделение сети продолжается, $B$ не знает о записи $x := 1$ на $A$.

Есть два варианта:
\begin{enumerate}
    \item $B$ отвечает старым значением. Тогда нарушается согласованность.
    \item $B$ не отвечает до восстановления связи. Тогда нарушается доступность.
\end{enumerate}

Следовательно, при сетевом разделении невозможно одновременно обеспечить и строгую согласованность, и доступность.
\end{proof}

\section{Практические рекомендации по разработке}

\subsection{Для программ над общей памятью}

\begin{enumerate}
    \item Минимизировать разделяемое изменяемое состояние.
    \item Использовать неизменяемые структуры данных, где возможно.
    \item Защищать все общие данные синхронизацией.
    \item Соблюдать единый порядок захвата блокировок.
    \item Не удерживать блокировки во время сетевых или длительных операций.
    \item Предпочитать высокоуровневые примитивы низкоуровневым.
    \item Тестировать программу под высокой нагрузкой.
    \item Использовать статические и динамические анализаторы гонок.
\end{enumerate}

\subsection{Для распределенных программ}

\begin{enumerate}
    \item Считать сеть ненадежной.
    \item Всегда задавать тайм-ауты.
    \item Делать операции идемпотентными, где возможно.
    \item Использовать уникальные идентификаторы запросов.
    \item Логировать важные переходы состояния.
    \item Проектировать повторную отправку сообщений.
    \item Разделять транспортные ошибки и ошибки бизнес-логики.
    \item Избегать предположения о глобальном времени.
    \item Проектировать систему с учетом частичных отказов.
\end{enumerate}

\section{Краткое сравнение моделей}

\begin{center}
\begin{tabular}{p{0.22\textwidth}p{0.32\textwidth}p{0.32\textwidth}}
\toprule
\textbf{Модель} & \textbf{Преимущества} & \textbf{Недостатки} \\
\midrule
Общая память & Быстрая коммуникация, удобство для многопоточных вычислений & Гонки, deadlock, сложная отладка \\
Передача сообщений & Явные зависимости, хорошо масштабируется на кластеры & Накладные расходы сети, сложность протоколов \\
RPC/RMI & Удобная абстракция удаленного вызова & Скрывает сетевые отказы и задержки \\
Акторы & Изоляция состояния, асинхронность & Сложность трассировки и согласования \\
MapReduce & Масштабируемая пакетная обработка & Не подходит для низколатентных задач \\
\bottomrule
\end{tabular}
\end{center}

\section{Итоговые тезисы}

\begin{enumerate}
    \item Параллельное программирование направлено прежде всего на ускорение вычислений, распределенное --- на масштабируемость, надежность и географическое распределение.
    \item Общая память удобна, но требует строгой синхронизации.
    \item Передача сообщений устраняет гонки по памяти, но вводит проблемы задержек, отказов и согласования.
    \item Логические часы позволяют рассуждать о причинности без глобального физического времени.
    \item Векторные часы позволяют точно определять причинную зависимость между событиями.
    \item Глобальное состояние распределенной системы должно учитывать не только состояния процессов, но и сообщения в каналах.
    \item Алгоритм Чанди--Лампорта строит согласованный снимок без остановки всей системы.
    \item Линеаризуемость --- важнейшее условие корректности параллельных объектов.
    \item Консенсус является фундаментальной задачей распределенных систем, но в полностью асинхронной модели с отказами имеет принципиальные ограничения.
    \item Эффективность параллельной программы ограничивается последовательной частью, синхронизацией и накладными расходами коммуникации.
\end{enumerate}

\end{document}
